\section{Coefficient normalization and the original Core-to-CoreBC implication}
\label{sec:bcn}

This chapter proves the full original implication
\(\mathrm{Core}\Rightarrow\mathrm{CoreBC}\), with arbitrary finite
chains of points.  The source is pinned at
\texttt{7cbf7a0c32fca1badf05931879a37c3f8a2b03b4}.
The assertion \(\mathrm{Core}\) is the antecedent throughout:
its patched support producer is not assumed to have been constructed
by this chapter.  In particular, none of the arguments below supplies
an initial \(\mathrm{TWAut}\).

The coefficient ring \(O\) is a complete discrete valuation ring,
with uniformizer \(\pi\), finite residue field \(k\) of characteristic
\(\ell\), and the original algebra map \(\mathbb Z_\ell\to O\).
An \(O\)-algebra \(R\) belongs to \(\mathrm{IsCO}(O,R)\) when it is
Noetherian local, complete for its maximal ideal, the map \(O\to R\)
is local, and it induces an isomorphism on residue fields.  These
are precisely the fields of the source record, including its
residue-surjectivity field.  No flatness or domain condition on
\(R\) is included in this notation.

For the extension \(O'\) constructed below, put
\[
 A_R=O'\otimes_O R,\qquad
 \mathscr T(R)=A_R[[s]],\qquad
 \iota_R(r)=1\otimes r,\qquad C_R:A_R\longrightarrow A_R[[s]].
 \tag{BCN.1}
\]
In the formula for \(\iota_R\), the tensor is regarded as a constant
series.  Thus \(\mathscr T(R)\), \(C_R\), \(\iota_R\), and \(s\) are
literally the source's \(\mathrm{Til}(O,O',R)\), \(\mathrm{Til.C}\),
\(\mathrm{Til.\iota}\), and \(\mathrm{Til.s}\).  They are not
completed tensor products in place of those objects.
For an \(O\)-algebra homomorphism \(f:R\to R'\), the map
\(\mathscr T(f)\) applies \(1\otimes f\) to each coefficient.

We write \(\Delta=u.D\) for the original determinant law, to
distinguish it from the integer \(D\) counting weight variables.
All identities concerning \(\Delta\) are identities of
multiplicative homogeneous polynomial laws after every commutative
coefficient extension.  They are not identities only of traces.

\subsection{The coefficient extension and finite-module calculus}

% BEGIN BCN1
\paragraph{BCN1. A single polynomial for all the prescribed finite fields.}\label{claim:BCN1}
Let \(K\) be any finite field supplied with an \(O\)-algebra structure,
and set \(Q=|K|\).  There is a monic polynomial \(P\in O[X]\), of
positive degree \(d\), whose reduction \(g\in k[X]\) is irreducible,
such that
\[
 k[X]/(g)\simeq K,\qquad
 P_{K'}\text{ splits over }K'
 \quad\hbox{for every finite \(O\)-field \(K'\) with }|K'|=Q .
 \tag{BCN.2}
\]
Here the second assertion uses the given \(O\)-algebra structure
on each \(K'\), without choosing an identification \(K'\simeq K\).

\emph{Proof.}
The kernel of \(O\to K\) is nonzero.  Otherwise the distinct powers
of the nonzero nonunit \(\pi\) would inject an infinite subset of
\(O\) into \(K\).  A nonzero prime of a discrete valuation ring is
its maximal ideal, so this map induces an embedding \(k\hookrightarrow K\).
Choose a generator \(\bar\alpha\) of the finite extension \(K/k\).
For completeness, the multiplicative group of a finite field is
cyclic: for every divisor of its exponent a polynomial of that
degree has at most that many roots; separating the prime-power
parts of the exponent produces an element of the full exponent,
and the root bound forces that exponent to equal the group order.
A generator of \(K^\times\) therefore generates \(K\) as a field.
Let \(g\) be its monic minimal polynomial and lift its nonleading
coefficients arbitrarily to \(O\), keeping the leading coefficient
equal to \(1\).  This defines \(P\).

Since \(\bar\alpha^Q=\bar\alpha\), polynomial division gives
\(g\mid X^Q-X\) in \(k[X]\).  Every map \(O\to K'\) as in the
statement again factors through an embedding of \(k\), and the
image of \(P\) is the corresponding image of \(g\).  It therefore
divides \(X^Q-X\) in \(K'[X]\).  The latter polynomial is the
product of \(X-c\) over all \(c\in K'\), since it is monic of
degree \(Q\) and vanishes at all \(Q\) elements.  Thus \(P_{K'}\)
splits.  This argument also covers differing embeddings of \(k\)
into the fields in the chain.
\hfill\(\square\)
% END BCN1

% BEGIN BCN2
\paragraph{BCN2. The original \(\mathrm{Oprime.oprime}\) output.}\label{claim:BCN2}
With \(P\) as in BCN1, let
\[
 O'=O[X]/(P),\qquad \alpha=X\bmod P .
 \tag{BCN.3}
\]
Then \(O'\) is a complete discrete valuation ring, finite free
of rank \(d\) over \(O\), and \(O\to O'\) is local.  Its maximal
ideal is \(\pi O'\), its residue field has cardinality \(Q\),
\(O[\alpha]=O'\), and \(P(\alpha)=0\).
With the composite \(\mathbb Z_\ell\to O\to O'\), it has the
original \(\mathrm{IsCoef}(\ell,O')\) structure.

\emph{Proof.}
Division by the monic \(P\) gives the \(O\)-basis
\(1,\alpha,\ldots,\alpha^{d-1}\).  In particular \(O\to O'\) is
injective and multiplication by \(\pi\) is injective.
Moreover
\[
 O'/\pi O'=k[X]/(g)\simeq K.                       \tag{BCN.4}
\]
The finite \(O\)-algebra \(O'\) is integral over \(O\).
The contraction of any maximal ideal of an integral extension
is maximal: after passage to the corresponding field quotient,
if the inverse of a nonzero element is integral, its monic
equation expresses that inverse in the base domain.
Every maximal ideal of \(O'\) therefore contains \(\pi O'\);
(BCN.4) shows that this is the unique maximal ideal.

The finite free coordinates also give
\(\bigcap_{r\ge0}\pi^rO'=0\) and completeness for these powers,
because these facts hold coordinatewise in \(O^d\).
For \(a\ne0\), choose the largest \(r\) with \(a\in\pi^rO'\)
and write \(a=\pi^r v\).  Then \(v\notin\pi O'\), so \(v\) is
a unit.  Products of two such expressions are nonzero, since
multiplication by every \(\pi^r\) is injective.  Thus \(O'\) is
a domain.  The least valuation of the nonzero elements of any
nonzero ideal gives a generator of that ideal.  Consequently
\(O'\) is a discrete valuation ring with uniformizer \(\pi\);
it is not a field because \(\pi\) is a nonzero nonunit.
Completeness and the finite residue characteristic in (BCN.4)
are exactly the three \(\mathrm{IsCoef}\) fields.  The generation
and polynomial assertions follow from (BCN.3), and BCN1 supplies
the final universally quantified splitting assertion of
\(\mathrm{oprime}\).
\hfill\(\square\)
% END BCN2

% BEGIN BCN3
\paragraph{BCN3. The local ring and the actual added variable.}\label{claim:BCN3}
For every \(R\) with \(\mathrm{IsCO}(O,R)\), the ring \(A_R\)
is finite free over \(R\), is local and complete, and satisfies
\[
 \mathfrak m_{A_R}=\mathfrak m_R A_R,\qquad
 A_R/\mathfrak m_R A_R\simeq k',\qquad
 k'=O'/\pi O'.
 \tag{BCN.5}
\]
The original \(\mathscr T(R)\) has \(\mathrm{IsCO}(O',\mathscr T(R))\);
\(\iota_R\) is local, and \(s\in\mathfrak m_{\mathscr T(R)}\).
In particular,
\[
 \mathfrak m_{\mathscr T(R)}=(\mathfrak m_{A_R},s).
 \tag{BCN.6}
\]

\emph{Proof.}
The basis from BCN2 identifies \(A_R\) with \(R^d\) as an
\(R\)-module.  The residue-field identification in
\(\mathrm{IsCO}(O,R)\) gives
\[
 A_R/\mathfrak m_R A_R
 =O'\otimes_O(R/\mathfrak m_R)\simeq O'/\pi O'=k'.
\]
Integrality over \(R\) forces every maximal ideal of \(A_R\)
to contract to \(\mathfrak m_R\), so this field quotient proves
(BCN.5).  A finite module over a Noetherian ring is Noetherian;
in particular all ideals of \(A_R\) are finitely generated.
Its \(\mathfrak m_{A_R}\)-adic completeness and separation hold
in its finite free \(R\)-coordinates.

We record the elementary power-series facts used here and later.
For a Noetherian ring \(A\), the ring \(A[[s]]\) is Noetherian.
One proof collects the possible leading coefficients at order
\(j\) of elements of a fixed ideal.  These ideals of \(A\) form
an increasing sequence, hence stabilize.  Choose finitely many
series realizing generators at the finitely many orders before
stabilization and at the stable order.  Successively subtract
multiples of these series to remove the leading term of any
element of the ideal.  At each step the order increases;
the coefficients of each multiplier converge as formal series.
This expresses the element as a finite sum of these chosen
generators with formal-series multipliers.  Iterating the
one-variable argument gives the same fact in finitely many
variables.

A series is a unit exactly when its constant term is a unit;
the inverse is obtained recursively.  Hence (BCN.6) and locality
follow.  A Cauchy sequence for \((\mathfrak m_{A_R},s)\) has
coefficientwise \(\mathfrak m_{A_R}\)-adic limits, and the defining
congruences show that its series of limits is the required limit.
The residue field is \(k'\), the map from \(O'\) induces the
identity on it, and \(\iota_R\) contracts the maximal ideal to
\(\mathfrak m_R\).  This proves every asserted \(\mathrm{IsCO}\)
field and the original \(\mathrm{til\_co}\) output for our
particular \(O'\).
\hfill\(\square\)
% END BCN3

% BEGIN BCN4
\paragraph{BCN4. Finite presentations, series, and exact module transport.}\label{claim:BCN4}
Let \(R\) be a Noetherian \(O\)-algebra and \(M\) a finite
\(R\)-module.  There is a natural additive isomorphism
\[
 \mathscr T(R)\otimes_R M
 \ \simeq\ (O'\otimes_O M)[[s]],\qquad
 \Big(\sum_{r\ge0}a_rs^r\Big)\otimes m
 \longmapsto \sum_{r\ge0}a_r m\,s^r .
 \tag{BCN.7}
\]
The notation on the right denotes coefficient sequences, with
the natural convolution module action.
If \(S\to R\) is a homomorphism of Noetherian \(O\)-algebras and
\(M\) is finite over \(S\), with its \(S\)-action obtained from
its \(R\)-action, it follows that
\[
 \mathscr T(S)\otimes_S M\simeq
 \mathscr T(R)\otimes_R M,                         \tag{BCN.8}
\]
semilinearly for \(\mathscr T(S)\to\mathscr T(R)\).
Finite ideal quotients commute with \(\mathscr T\); in particular,
for every finitely generated ideal \(I\subset R\),
\[
 \mathscr T(R/I)\simeq
 \mathscr T(R)/(C_R(1\otimes I)),                  \tag{BCN.9}
\]
and this ideal consists exactly of the series whose coefficients
belong to \(O'\otimes_O I\).

\emph{Proof.}
Choose a finite presentation \(R^a\to R^b\to M\to0\).
Tensoring first with the finite free \(O'\) is exact.  Surjectivity
onto coefficient sequences follows by choosing a lift of each
coefficient in \((O'\otimes R)^b\).  The kernel consists of
sequences in the image of the fixed relation matrix; choose
a preimage for each coefficient and assemble the finitely many
coordinate series.  Thus this kernel is the image of the same
matrix over \(\mathscr T(R)\), proving (BCN.7).
The finite generation of the relation module is precisely where
Noetherianity is used.

For (BCN.8), \(M\) is also finite over \(R\), using the same
finite \(S\)-generators.  Both sides of (BCN.8) identify with the
right side of (BCN.7), and on a pure tensor the map is
\(a\otimes m\mapsto\mathscr T(f)(a)\otimes m\).
Its scalar law follows coefficient by coefficient, and hence on
finite sums.  For (BCN.9), finite generators of \(I\) express
every coefficient in the kernel with one fixed finite list;
assembling their coefficients proves both the kernel assertion
and surjectivity.

We will also use the following completeness consequence, with
its elementary justification explicit.  A finite module over
a complete Noetherian local ring is complete and separated,
and every submodule is closed.  To see the separation, the
Rees-module argument gives Artin--Rees: for a submodule \(N\)
of a finite module \(M\), the graded module
\(\bigoplus_j(\mathfrak m^jM\cap N)\) is a submodule of the
finite Rees module over the Noetherian ring
\(R[\mathfrak m t]\); finite homogeneous generators imply
\(\mathfrak m^jM\cap N\subset\mathfrak m^{j-c}N\) for one \(c\)
and all \(j\ge c\).  Applying this to the intersection of all
powers, which is itself finite, gives \(N=\mathfrak m N\);
the determinant proof of Nakayama gives \(N=0\).
Thus submodules are closed by applying separation to the
quotient.  For completeness, lift a compatible Cauchy sequence
to a finite free presentation and adjust successive lifts by
relations so that their differences lie in successive powers
of \(\mathfrak m\).  The free coordinates converge; their image
is the required limit.  This proves the consequence without
assuming completeness of an arbitrary infinite tensor product.
\hfill\(\square\)
% END BCN4

% BEGIN BCN5
\paragraph{BCN5. The exact weight, auxiliary, and splitting algebras.}\label{claim:BCN5}
Set
\[
 \Lambda=O[[z_k:k\in\mathrm{Fin}(D)]],\qquad
 \Sigma=\Lambda[[y_{ji}:(j,i)\in\mathrm{Fin}(q)\times
 \mathrm{Fin}(n)]].
\]
The empty summand in the source's old variable type is understood.
The finite free coefficient basis and interchange of coefficient
indices give natural isomorphisms
\[
 \begin{aligned}
 \tau:\mathscr T(\Lambda)&\simeq
 \Lambda'=O'[[s,z_k]],\\
 \tau_\Sigma:\mathscr T(\Sigma)&\simeq
 \Sigma'=\Lambda'[[y_{ji}]] .
 \end{aligned}                                                   \tag{BCN.10}
\]
They send the old variables to the identically indexed new
variables, the added \(s\) to the literal variable indexed by
\(\mathrm{inl}(())\), and old coefficients by \(O\to O'\).
They commute with the coefficient inclusion
\(\Lambda\to\Sigma\) and with the constant coefficient in all
\(y\)-variables.  They carry the original ideal
\[
 \delta_a=\big((1+y_{ji})^{\ell^{a_j}}-1\big)_{j,i}               \tag{BCN.11}
\]
to the identically defined \(\delta'_a\).
For the original splitting algebra
\[
 \mathrm{Sp}(P_j,u^0_{ji})
 =B[[Y_{ji}]]/
 \left(\operatorname{coeff}_{X^r}
   \left[\prod_i(X-u^0_{ji}-Y_{ji})-P_j(X)\right]\right)_{j,r},
                                                               \tag{BCN.12}
\]
there is an isomorphism
\[
 \tau_{\mathrm{Sp}}:\mathscr T(\mathrm{Sp}(P,u^0))
 \simeq \mathrm{Sp}(P^\iota,\iota u^0)                         \tag{BCN.13}
\]
over \(\mathscr T(B)\).  It carries each original distinguished
root \(u^0_{ji}+Y_{ji}\) to the corresponding new distinguished root.

\emph{Proof.}
An element in the first ring of (BCN.10) is an array with an
\(s\)-index, a \(z\)-multi-index, and one of the finitely many
\(O\)-basis indices of \(O'\).  Interchanging the first two
indices gives exactly the indicated \(O'\)-coefficient series.
Addition and multiplication agree because each coefficient of
a product is a finite sum; no continuity assertion about an
arbitrary ring homomorphism is needed.  The same argument adds
the \(y\)-indices and proves every stated coefficient and
variable identity.  Equation (BCN.11) is a finite list of
polynomials, so its ideal is carried exactly to \(\delta'_a\).

For (BCN.13), first interchange the \(s\)-index and all \(Y\)-indices.
The ideal in (BCN.12) is finitely generated: there are finitely
many \(j\), and only finitely many nonzero coefficients in
each displayed polynomial.  Equation (BCN.9) therefore
commutes with this quotient.  Its generators map to the
identical new coefficient relations.  The image of a root
and the \(\mathscr T(B)\)-algebra identity follow directly on
coefficients.  These are the original \(\mathrm{til\_lam}\),
\(\mathrm{sinf\_iso}\), and \(\mathrm{til\_sp}\) compatibility
identities needed below.
\hfill\(\square\)
% END BCN5

% BEGIN BCN6
\paragraph{BCN6. Finiteness of the actual splitting module.}\label{claim:BCN6}
Let \(B\) be complete Noetherian local.  The algebra in
(BCN.12) is finite over \(B\), including when some prescribed
residue roots coincide.  It is therefore Noetherian.
If \(M\) is finite over \(\Lambda\), its \(B\)-action agrees
with \(\lambda:\Lambda\to B\), and
\(\theta:R\to S=\mathrm{Sp}(P,u^0)\) has
\[
 \theta(s_R(C_\Sigma a))=\lambda(a)\quad(a\in\Lambda),
                                                               \tag{BCN.14}
\]
where the right side is mapped to \(S\), then \(S\otimes_B M\)
is finite as an \(R\)-module through \(\theta\).

\emph{Proof.}
Write \(S\) for (BCN.12) and reduce modulo \(\mathfrak m_B\).
If the quotient is zero there is nothing to check at this step.
Otherwise \(k\) embeds in it, and the polynomial identity in
(BCN.12) shows that the image of \(P_j\) is nonzero: its
leading coefficient at degree \(n\) is \(1\).  Each variable
\(Y_{ji}\) lies in the Jacobson radical of the quotient of
\(k[[Y]]\), because \(1+Y_{ji}v\) is a unit for every \(v\).
It satisfies the nonzero polynomial \(P_j(u^0_{ji}+X)\) over
\(k\).  Factor that polynomial as \(X^a b(X)\), with \(b(0)\ne0\).
Then \(b(Y_{ji})\) is a unit, so \(Y_{ji}^a=0\).
Finitely many such nilpotent variables generate a nilpotent
ideal: an exponent exceeding the sum of their exponents
forces every monomial to contain a vanishing power.
It follows that \(S/\mathfrak m_BS\) is finite dimensional.

The ring \(S\) is a quotient of the complete Noetherian local
ring \(B[[Y]]\), and is complete for its maximal ideal by
BCN4.  Nilpotence of the residual variables gives
\(\mathfrak m_S^c\subset\mathfrak m_BS\) for some \(c\);
the \(\mathfrak m_S\)-adic and \(\mathfrak m_B\)-adic
topologies are therefore equivalent.  Lift a finite basis
of \(S/\mathfrak m_BS\).  Successively express the error in
\(\mathfrak m_B^rS\) using these same lifts, and let the
coefficients converge in \(B\).  Completeness and separation
show that the lifts span \(S\) over \(B\).
This also deals with a zero residue quotient: the same
iteration gives \(S=0\).

Choose finite \(B\)-generators \(s_j\) of \(S\) and finite
\(\Lambda\)-generators \(m_k\) of \(M\).
Every tensor is a sum of \(s_j\otimes m\), after moving its
\(B\)-coefficients to \(M\); express each \(m\) in the \(m_k\).
Equation (BCN.14) moves each resulting \(\Lambda\)-coefficient
to the \(R\)-action.  Thus the finite set \(s_j\otimes m_k\)
spans \(S\otimes_B M\) over \(R\), exactly as required by
\(\mathrm{exists\_finite\_spanning\_taylorWiles\_tensor}\).
\hfill\(\square\)
% END BCN6

\subsection{The entire original automorphic packet}

% BEGIN BCN7
\paragraph{BCN7. Residual data, full pairs, and dense generation.}\label{claim:BCN7}
Let \(r\) be the original residual datum over \(O\).
Push its residual characters through \(k\hookrightarrow k'\)
and its \(O\)-valued characters through \(O\to O'\).
This gives a residual datum \(r'\), with the exact \(\mathrm{RBC}\)
relations and
\[
 r'.I_p(v)=r.I_p(v),\qquad
 r'.\bar D=r.\bar D\otimes_k k',\qquad
 r'.\bar\chi_i(g)=1\ \Longleftrightarrow\ r.\bar\chi_i(g)=1 .
                                                               \tag{BCN.15}
\]
For \(u=(\Delta,\psi)\) over \(B\), define
\[
 u'=(\Delta\otimes_B\mathscr T(B),\,\iota_B\psi).
                                                               \tag{BCN.16}
\]
It is a full original \(\mathrm{Pair}\) over \(\mathscr T(B)\).
It has the original \(\mathrm{Maps}\) identity under \(\iota_B\).
Adaptedness to the same \(T\) is preserved, and
\[
 \mathrm{GenBy}_O(\mathrm{pairVals}(u))
 \ \Longrightarrow\
 \mathrm{GenBy}_{O'}
       (\mathrm{pairVals}(u')\cup\{s\}).                         \tag{BCN.17}
\]

\emph{Proof.}
The residue-field injection preserves every character kernel,
every equality to \(1\), and the pairwise distinction of the
original residual characters.  It preserves the open kernels
and all commutator identities.  The cyclotomic coefficient
map is the same composite \(\mathbb Z_\ell\to O\to O'\);
therefore the polarized involution \(\star\) commutes with
coefficient extension.  Evaluating the defining polarization
identity after any further commutative algebra gives the
required residual polarization.  This proves every residual
field and (BCN.15).

For the pair, naturality gives
\[
 \chi_{\Delta\otimes\iota}(\iota y)
      =\iota(\chi_\Delta(y)),\qquad
 \Lambda_j^{\Delta\otimes\iota}(\iota y)
      =\iota(\Lambda_j^\Delta(y)).                              \tag{BCN.18}
\]
The local map \(\iota\) sends \(\mathfrak m_B^r\) into
\(\mathfrak m_{\mathscr T(B)}^r\), so all the determinant and
character continuity conditions follow at the same adic
depth.  The residual square commutes by (BCN.5).
The \(\mu\)-values of the local characters commute with
\(O\to O'\); the polarization, ordered ordinary products,
and inertia trace identities are the images of the old
identities, by (BCN.18) and the equality of \(I_p\)'s.
This checks the full \(\mathrm{Pair}\), including
\(\mathrm{P}\), \(\mathrm{Ord}\), and \(\mathrm{I}\).
Adaptedness now follows immediately from (BCN.15).

Finally fix a depth \(r\).  Truncate a series in
\(\mathscr T(B)\) modulo \(s^r\).  Each of its finitely many
tensor coefficients is a finite sum of \(O'\)-multiples of
elements of \(B\).  Approximate each such element modulo
\(\mathfrak m_B^r\) by the \(O\)-algebra generated by the
original pair values.  The images of these generators are
pair values of \(u'\), by (BCN.18) and the definition of
the new characters.  Since \(s\) is in the new maximal
ideal, both the truncated tail and all coefficient errors
belong to its \(r\)-th power.  This is exactly the quantified
adic-density definition of \(\mathrm{GenBy}\), proving
(BCN.17).
\hfill\(\square\)
% END BCN7

% BEGIN BCN8
\paragraph{BCN8. The base module, its free structure, and its support.}\label{claim:BCN8}
For an original \(\mathrm{TWBase}(u,J,\mathrm{Empty},D)\),
define
\[
 \begin{aligned}
 J'&=J\,\mathscr T(B),\\
 \lambda'&=\mathscr T(\lambda)\tau^{-1}:\Lambda'\to\mathscr T(B),\\
 M'&=\mathscr T(B)\otimes_B M .
 \end{aligned}                                                   \tag{BCN.19}
\]
Keep the original indices \(v_k,i_k,\gamma_k\), and replace each
weight coefficient \(c_k\) by its image in \(O'\).
These choices give the full original
\(\mathrm{TWBase}(u',J',\mathrm{Unit},D)\), with
\[
 \lambda'(s)=s,\qquad
 M'\simeq\Lambda'\otimes_\Lambda M
 \quad\hbox{as \(\Lambda'\)-modules}.                           \tag{BCN.20}
\]
The displayed module is therefore free of the same finite rank.

\emph{Proof.}
The weight residue identity is its image in \(k'\).
The weight-variable identity follows from (BCN.10),
(BCN.16), and the original
\(\lambda(z_k)=c_k\psi_{v_k,i_k}(\gamma_k)-1\).
The \(\Lambda\)-free finite module \(M\) is finite over \(B\):
its \(\Lambda\)-basis spans it over \(B\), since
\(a m=\lambda(a)m\).  Equation (BCN.8), followed by \(\tau\),
proves (BCN.20).  Define the new \(\Lambda'\)-action through
\(\lambda'\); the required scalar identity is then literal.

It remains to check the support field for every prime
\(\mathfrak p'\subset\mathscr T(B)\), rather than only at
the selected points.  Put \(\mathfrak p=\iota_B^{-1}\mathfrak p'\).
The residue fiber is
\[
 \kappa(\mathfrak p')\otimes_B M
 \simeq\kappa(\mathfrak p')\otimes_{\kappa(\mathfrak p)}
       (M_{\mathfrak p}/\mathfrak pM_{\mathfrak p}).
\]
A field extension is faithful on vector spaces.  Nakayama
for the finite modules at these local rings therefore says
that \(M'_{\mathfrak p'}\ne0\) exactly when \(M_{\mathfrak p}\ne0\).
For a finite module this is equivalent to containment of its
annihilator in the prime: if its localization vanishes, finitely
many generators have a common annihilator outside the prime;
the converse is immediate.  The old support field gives
\[
 \operatorname{Ann}_{\mathscr T(B)}M'\subset\mathfrak p'
 \ \Longleftrightarrow\
 J\subset\mathfrak p
 \ \Longleftrightarrow\
 J'\subset\mathfrak p'.                                       \tag{BCN.21}
\]
This is exactly the new support field; no equality between
unextended and extended annihilator ideals has been assumed.
The first identity in (BCN.20) follows from the literal
\(s\)-variable in (BCN.10).
\hfill\(\square\)
% END BCN8

% BEGIN BCN9
\paragraph{BCN9. Every auxiliary ring and its full module structure.}\label{claim:BCN9}
Fix any level \(N\) of the original \(\mathrm{TWAut}\).
Write \(R_N,p_N,s_N,U_{N,ji},M_N,a_{N,j},r_N\) for its
\(\mathrm{TWRing}\) fields, using \(r_N\) here for its module
rank.  Define
\[
 \begin{aligned}
 R_N'&=\mathscr T(R_N),&
 p_N'&=p_N\otimes\iota_{R_N},\\
 s_N'&=\mathscr T(s_N)\tau_\Sigma^{-1},&
 U'_{N,ji}&=\iota_{R_N}(U_{N,ji}),\\
 M_N'&=\mathscr T(R_N)\otimes_{R_N}M_N .
 \end{aligned}                                                   \tag{BCN.22}
\]
Keep the same exponents \(a_{N,j}\), the same rank \(r_N\),
the same group \(T.G_N\), and the same homomorphism \(T.h_N\).
Push the residual datum as in BCN7.  These data satisfy every
field of the original new \(\mathrm{TWRing}\).

\emph{Proof.}
BCN3 gives the required \(\mathrm{IsCO}\), and BCN7 gives
the residual compatibility and full pair.
The weight relation follows coefficientwise from (BCN.10)
and the old relation.  The condition
\(U_{N,ji}-u^0_{ji}\in\mathfrak m_{R_N}\) maps to the new
maximal ideal, so the residual-root condition is unchanged.
The numerical inequalities \(N\le a_{N,j}\) are identical.

The given additive equivalence
\(M_N\simeq(\Sigma/\delta_{a_N})^{r_N}\) and its scalar law
make \(M_N\) a finite \(\Sigma\)-module, hence a finite
\(R_N\)-module.  Apply (BCN.8), then (BCN.9) and (BCN.10):
\[
 \mathscr T(R_N)\otimes_{R_N}M_N
 \simeq \mathscr T(\Sigma)\otimes_\Sigma M_N
 \simeq
 \big(\Sigma'/\delta'_{a_N}\big)^{r_N}.                          \tag{BCN.23}
\]
This supplies the exact new additive equivalence and its
\(\Sigma'\)-scalar identity.  Thus the construction preserves
the original module presentation, rather than substituting
only its rank or its support.
\hfill\(\square\)
% END BCN9

% BEGIN BCN10
\paragraph{BCN10. Both comparison maps and all Taylor--Wiles relations.}\label{claim:BCN10}
At the same level \(N\), let \(S=\mathrm{SpAt}(u,\ldots)\)
be the original splitting algebra, \(X=M_N/(yM_N)\), and
\(X_0=S\otimes_B M\).
Let
\(\theta:R_N\to S\), \(\phi:X_0\to X\), and
\(\psi:X\to X_0\) be the original \(\mathrm{TWMaps}\).
Using \(\tau_{\mathrm{Sp}}\), define
\[
 \theta'=\tau_{\mathrm{Sp}}\mathscr T(\theta).
                                                               \tag{BCN.24}
\]
There are canonical identifications
\[
 X'=\mathscr T(R_N)\otimes_{R_N}X,\qquad
 X_0'\simeq\mathscr T(S)\otimes_S X_0
       \simeq\mathscr T(R_N)\otimes_{R_N}X_0 .                  \tag{BCN.25}
\]
Tensor \(\phi,\psi\) over \(R_N\) and transport through
(BCN.25).  The resulting \(\phi',\psi'\), together with
\(\theta'\), satisfy every original \(\mathrm{TWMaps}\) field,
with the same exponent \(\mathrm{ex}\), the same
\(Nq,\mathrm{ph},\mathrm{ta}\), and the same ordered products.

\emph{Proof.}
The first identification in (BCN.25) is the right-exact
tensor-quotient identity; the new ideal is generated by
the images of exactly the old \(s_N(y_{ji})\).
Associativity of tensor products and BCN8 give the first
identification for \(X_0'\).  To justify the second one,
BCN6 applies to the original coefficient identity
\(\theta s_N(Ca)=\lambda(a)\).
Thus \(X_0\) is finite over \(R_N\) through \(\theta\);
it is finite over \(S\) as well.  Both rings are Noetherian,
so BCN4 gives the required semilinear identification.
The original scalar laws say exactly that \(\phi,\psi\)
are \(R_N\)-linear for this action.  Tensoring them therefore
gives the new scalar laws.

Let \(F=fEl(Nq,U_N)\).  It is the literal finite product
\[
 F=\prod_{j,i}\ \prod_{i'\ne i}
       (U_{N,ji}-U_{N,ji'})(U_{N,ji}-Nq_j U_{N,ji'}).
                                                               \tag{BCN.26}
\]
Its new value is \(\iota_{R_N}(F)\).
The two original composition identities tensor to
\(\phi'\psi'=\iota(F)\) on \(X'\) and
\(\psi'\phi'=\theta'(\iota(F))\) on \(X_0'\).

The full determinant and all local-character identities
for \(\theta'\) follow by functoriality from the commuting
square \(\theta'\iota=\iota_S\theta\).
BCN5 makes the coefficient identity commute with the
constant coefficient in the \(y\)'s, for every series,
not only for polynomials.  It also gives the distinguished
root identity \(\theta'(U'_{N,ji})=u'_{ji}\).

For clarity, the remaining annihilator conditions involve
the original noncommutative ordered expressions
\[
 \begin{aligned}
 q_1&=\prod_{i\text{ in order}}([\mathrm{ph}_j]-U_{N,ji}),\\
 q_2&=\left(\prod_{k\text{ in order},\,k\ne i}
             ([\mathrm{ph}_j]-U_{N,jk})\right)
             ([\mathrm{ta}_j]-1-s_N(y_{ji})).
 \end{aligned}                                                   \tag{BCN.27}
\]
Coefficient extension preserves their displayed order.
It also preserves the trace of
\([g]q_1\) and \([g]q_2\), by (BCN.18).
Hence each new annihilating scalar
\((F\operatorname{tr}([g]q_a))^{\mathrm{ex}}\)
is the image of the old one.  If a scalar kills \(M_N\),
its image kills every pure tensor in \(M_N'\), and therefore
the whole module.  This verifies both \(\mathrm{rel1}\) and
\(\mathrm{rel2}\), with no change of exponent or level.
These checks exhaust the fields of \(\mathrm{TWMaps}\).
\hfill\(\square\)
% END BCN10

% BEGIN BCN11
\paragraph{BCN11. The literal Unit-parameter \(\mathrm{TWAut}\).}\label{claim:BCN11}
Every supplied original packet
\[
 \mathcal A:\mathrm{TWAut}(u,J,\mathrm{Empty},D,T)
\]
gives, by BCN8--BCN10, an actual
\[
 \mathcal A':
 \mathrm{TWAut}(u',J',\mathrm{Unit},D,T),\qquad
 \mathcal A'.b.\lambda(X_{\mathrm{inl}(())})=s .
                                                               \tag{BCN.28}
\]
The new residue-root lifts are the images in \(O'\) of the
old lifts.  All levels \(N\), all exponent arrays, all ranks,
and the one common \(\mathrm{ex}\) are preserved.

\emph{Proof.}
Take the base packet from BCN8, and at every \(N\) take the
ring packet of BCN9 and the maps of BCN10.
The residue equation for each root lift is the image under
\(k\hookrightarrow k'\) of the original equation, using
the residual characters of BCN7.  These are all fields
of \(\mathrm{TWAut}\).  The last equality is (BCN.20) with
the original variable label restored.
The same group packet \(T\), not an isomorphic replacement,
indexes the entire construction.
\hfill\(\square\)
% END BCN11

\subsection{Lifting the same points and their entire prime chain}

% BEGIN BCN12
\paragraph{BCN12. A prescribed prime lifts to the same series field.}\label{claim:BCN12}
Let \(K_i\) be any of the finite \(O\)-fields with
\(|K_i|=Q\), let \(f_i:B\to K_i[[t_i]]\) be its original
\(O\)-algebra map, and put \(p_i=\ker f_i\).
For every prime \(P_i\subset A_B\) contracting to \(p_i\),
there are an \(O'\)-algebra structure on the same \(K_i\)
and an \(O'\)-algebra map
\[
 f_i':\mathscr T(B)\longrightarrow K_i[[t_i]]
\]
such that
\[
 O'\longrightarrow K_i\text{ is onto},\qquad
 f_i'\iota_B=f_i,\qquad f_i'(s)=t_i,\qquad
 \ker(f_i' C_B)=P_i.                                           \tag{BCN.29}
\]
The coefficient extension \(O'\) is the one common ring
chosen in BCN2 for the whole chain.

\emph{Proof.}
The presentation (BCN.3) gives
\(A_B=B[X]/(P)\).  Set \(R=B/p_i\), a subring of \(K_i[[t_i]]\),
and \(F=\operatorname{Frac}(R)\subset K_i((t_i))\).
The prime \(P_i\) is disjoint from \(B\setminus p_i\).
After quotienting by \(p_i\) and localizing, it gives a
prime of \(F[X]/(P)\), hence an irreducible factor \(b(X)\)
of \(P\) over \(F\).
By BCN1, \(P\) splits as a product of linear factors with
roots in the constant field \(K_i\).
The factor \(b\), after extension to \(K_i((t_i))\), thus
has a root \(c\in K_i\).
Evaluation at \(c\) embeds \(F[X]/(b)\) into \(K_i((t_i))\).
Restricting this embedding gives
\[
 x:A_B\longrightarrow K_i[[t_i]],\qquad
 x(1\otimes b)=f_i(b),\quad x(\alpha\otimes1)=c,\quad
 \ker x=P_i .                                                 \tag{BCN.30}
\]
The image lies in the series ring, because \(A_B\) is generated
over \(B\) by \(\alpha\).  Equality of the kernel follows
after localization and then by contraction, since \(P_i\)
was disjoint from the localized multiplicative set.

Every element of \(O'\) is a polynomial in \(\alpha\) with
coefficients in \(O\), whose images under \(f_i\) are the
given constants.  Thus \(x|_{O'}\) is constant and determines
a map \(\varphi_i:O'\to K_i\), extending the original map
from \(O\).  It factors through \(k'\), by the finite-field
kernel argument of BCN1 applied to \(O'\).  The resulting
embedding \(k'\hookrightarrow K_i\) is onto because both
fields have cardinality \(Q\).  Give \(K_i\) this
\(O'\)-algebra structure.

For \(a(s)=\sum_{r\ge0}a_rs^r\), define
\[
 f_i'(a)=\sum_{r\ge0}x(a_r)t_i^r .                              \tag{BCN.31}
\]
This is a well-defined series: each coefficient involves
only finitely many terms \(r\).  The same finite convolution
check proves additivity, multiplicativity and preservation
of \(1\); (BCN.30) proves the \(O'\)-algebra law.
Formula (BCN.31) restricts to \(x\) on constants and sends
\(s\) to \(t_i\).  This proves all of (BCN.29), including
the prescribed kernel on the coefficient tensor ring.
\hfill\(\square\)
% END BCN12

% BEGIN BCN13
\paragraph{BCN13. The exact \(\mathrm{Pt}\), \(\mathrm{Kap}\),
and Selmer bounds do not change.}\label{claim:BCN13}
Suppose the original point is
\[
 x_i:\mathrm{Pt}\bigl(T,u,D+nT.q\bigr).
\]
Using BCN12, define \(x_i'\) with the same
\[
 x_i.x,\quad x_i.k,\quad x_i.al,\quad
 x_i.eig,\quad x_i.b1,\quad x_i.b2,
\]
and with coefficient algebra and map \(f_i'\) from (BCN.29).
Then \(x_i'\) is a point of
\(\mathrm{Pt}(T,u',D+nT.q)\), and it satisfies the
normalization required by \(\mathrm{Core}\).

\emph{Proof.}
The \(\mathrm{NP}\) record is literally unchanged: the field,
series ring, representation, polarization form, block size,
span condition, local flags and character data are the old ones.
The \(\mathrm{Kap}\) record and its cocycles are also unchanged.
The eigenvalues and their distinctness and characteristic
polynomial identities therefore remain identical.

The two bounds do not mention the coefficient ring \(O\).
For \(\mathrm{B1}\), keep its original witness \(C\) in
\[
 \exists C\ \forall m\ \forall^{\mathrm{eventual}}N:
 \quad
 |\mathrm{Sel1}_{N,m}|
 \le |\mathrm{Bd1}_{N,m}|\,|K_i|^{\,m(D+nT.q)+C},
                                                               \tag{BCN.32}
\]
including its finiteness assertion.  All sets and maps in
this formula are unchanged.
For \(\mathrm{B2}\), keep its original \(C\), its order
\(\forall m\,\forall^{\mathrm{eventual}}N\), its universal
array of cocycles, its coefficient vector with a unit
entry, and its \(\mathrm{Sel2}\) witness.  Every equality
continues to use the same \(t_i^C\) and the same
\(\mathrm{Kap}\).  Thus no larger constants or later
eventual thresholds are introduced.

Finally the full-law identity is
\[
 (u'.D)\otimes f_i'
  =\Delta\otimes(f_i'\iota_B)
  =\Delta\otimes f_i
  =\det(x_i.x.\rho),
\]
and the local characters obey the same identity.
This checks \(\mathrm{Pt.maps}\), rather than only its
trace consequence.  Surjectivity \(O'\to K_i\) and
\(f_i'(\mathcal A'.b.\lambda(X_{\mathrm{inl}(())}))=t_i\)
are BCN12 and BCN11.
In particular the integers \(n\), \(x_i.x.a\), and \(x_i.k.g\)
are preserved exactly.
\hfill\(\square\)
% END BCN13

% BEGIN BCN14
\paragraph{BCN14. Simultaneous lifting of every finite prime chain.}\label{claim:BCN14}
The map \(B\to A_B\) is finite free of positive rank,
hence faithfully flat and integral.
Given primes
\[
 p_0,\ldots,p_h,\qquad
 e_j\subset p_j\cap p_{j+1}\quad(0\le j<h),
\]
there are primes \(P_i,E_j\subset A_B\), contracting to
\(p_i,e_j\), such that
\[
 E_j\subset P_j\cap P_{j+1}.                                   \tag{BCN.33}
\]
Moreover every prime \(q\subset P\cap B\) has a prime
\(Q\subset P\) contracting to \(q\).

\emph{Proof.}
The \(R=B\) instance of the basis in BCN2 gives a free
module of positive rank, so tensoring with \(A_B\) is
exact and faithful.  Finite algebra extensions are integral.

Here are the prime-lifting facts in the needed orientations.
For a flat local map of local rings, faithful flatness
follows directly: a nonzero cyclic submodule \(R/I\) of
any nonzero module stays nonzero after tensoring, since
\(I\) lies in the maximal ideal and its extension is proper;
flatness preserves its injection.  Localizing \(B\to A_B\)
at \(P\) therefore gives a faithfully flat map
\(B_{P\cap B}\to(A_B)_P\).
For \(q\subset P\cap B\), tensor this map with
\(\operatorname{Frac}(B_{P\cap B}/q)\).
Its fiber is nonzero.  A prime in this fiber, contracted
back to \(A_B\), is a prime \(Q\subset P\) over \(q\).
This proves the asserted going-down statement.
The same fiber argument for the original faithfully flat
map gives an initial prime over any \(p_0\).

For going up, quotient an integral extension by a prime
over \(e_j\), and localize the remaining base at \(p_{j+1}\).
Choose a maximal ideal in the resulting nonzero integral
algebra.  Its contraction is the maximal ideal of the
local base, by the integral field argument in BCN2.
Its inverse image is a prime over \(p_{j+1}\) containing
the chosen prime over \(e_j\).

Start with any \(P_0\) over \(p_0\).  Going down gives
\(E_0\subset P_0\) over \(e_0\); going up gives
\(P_1\supset E_0\) over \(p_1\).  Repeat for the finite
chain.  This constructs (BCN.33), including \(h=0\),
when there are no intermediate primes.
\hfill\(\square\)
% END BCN14

% BEGIN BCN15
\paragraph{BCN15. Prime extension to the actual series ring.}\label{claim:BCN15}
For a prime \(Q\subset A_B\), put
\(\widetilde Q=Q\,\mathscr T(B)\), generated by its
constant series.  Then
\[
 \mathscr T(B)/\widetilde Q\simeq(A_B/Q)[[s]],\qquad
 \widetilde Q\text{ is prime},\qquad
 \iota_B^{-1}\widetilde Q=Q\cap B .                             \tag{BCN.34}
\]
If \(Q\subset P_i\), then
\(\widetilde Q\subset\ker f_i'\) for the lift in BCN12.

\emph{Proof.}
The ring \(A_B\) is Noetherian, so \(Q\) has finitely many
generators.  Coefficientwise quotienting a series is
surjective, and assembling coefficients in these fixed
generators shows that its kernel is precisely
\(Q\mathscr T(B)\), as in BCN4.
A formal-series ring over a domain is a domain: the
lowest nonzero coefficients of two series have nonzero
product.  This proves primality and contraction on
constant series.  The last assertion follows because
\(f_i'\) kills \(C_B(P_i)\), hence kills the ideal
generated by \(C_B(Q)\).
\hfill\(\square\)
% END BCN15

\subsection{Closed relations and nonsemisimplicity without an extra
finite-generation premise}

% BEGIN BCN16
\paragraph{BCN16. Polynomial-law Cayley--Hamilton generators commute
with arbitrary coefficient extension.}\label{claim:BCN16}
For every commutative ring map \(B\to T_0\) and every full
determinant law \(\Delta\) on \(B[G]\),
\[
 \big\langle\mathrm{CH}(\Delta\otimes_B T_0)\big\rangle
 =
 \big\langle\operatorname{image}(\mathrm{CH}(\Delta))\big\rangle
 \quad\hbox{in }T_0[G].                                        \tag{BCN.35}
\]
The brackets denote the ordinary two-sided ideal.
There is no assumption that \(n!\) is invertible.

\emph{Proof.}
The source's \(\mathrm{CH}\) set consists of all coefficients,
in independent commuting variables \(Z_1,\ldots,Z_m\), of
\[
 \chi_\Delta\!\left(\sum_j Z_jy_j\right)
                \!\left(\sum_j Z_jy_j\right),
 \qquad y_j\in B[G].                                           \tag{BCN.36}
\]
The characteristic polynomial in this formula is evaluated
at the same group-algebra element; its coefficients are
central scalars.  Formula (BCN.36) is a finite polynomial
in the \(Z_j\)'s with finite group support.

For a new list \(y_j\in T_0[G]\), let \(S\) be the finite
union of its supports and write
\(y_j=\sum_{g\in S}c_{jg}[g]\).
Form (BCN.36) over \(B\), now using the elements \([g]\)
and independent variables \(W_g\).  Naturality of the full
law identifies the new expression with its coefficient
extension followed by the substitution
\(W_g=\sum_j c_{jg}Z_j\).
Every coefficient of this substitution is a finite
\(T_0\)-linear combination of the images of the old
coefficients.  This proves one containment of ideals.
The reverse containment follows by using lists of
old elements and the same variables after extension.
All coefficient comparisons are in polynomial rings,
so no polarization division, trace reconstruction,
or factorial cancellation occurs.
\hfill\(\square\)
% END BCN16

% BEGIN BCN17
\paragraph{BCN17. Closed group-algebra ideals under finite free
coefficient extension with cofinal local topologies.}\label{claim:BCN17}
More generally in this paragraph, let \(B\to A\) be a finite
free local homomorphism of local rings, and suppose that,
for one positive integer \(c\),
\[
 \mathfrak m_B A\subset\mathfrak m_A,\qquad
 \mathfrak m_A^c\subset\mathfrak m_B A.                         \tag{BCN.37a}
\]
The residue fields need not be equal.
For a topological group \(G\), a set \(Z\subset B[G]\), and
a determinant law \(\Delta\), write
\[
 I_B=
 \bigcap_{N\ge0,\,H\triangleleft_{\rm open}G}
 \left\langle
  \mathrm{CH}(\Delta),\,Z,\,\mathfrak m_B^N[1],\,
  [h]-1\ (h\in H)\right\rangle .                              \tag{BCN.37}
\]
This is exactly \(\Delta.\mathrm{chBar}(Z)\).
Under \(B\to A\), with the image set \(Z_A\) and
the extended determinant, the corresponding closed ideal
is exactly
\[
 I_A=A\,I_B\subset A[G].                                        \tag{BCN.38}
\]

\emph{Proof.}
Fix a finite \(B\)-basis of \(A\).
For a fixed pair \(N,H\), BCN16 shows that the ordinary
ideal with new Cayley--Hamilton and \(Z\)-relations,
the same \(H\)-relations, and the scalar ideal
\(\mathfrak m_B^N A\), is the extension of the old
ordinary ideal at \(N,H\).
A two-sided ideal of \(B[G]\), extended by central
coefficients, is just its \(A\)-linear span: multiplication
by group elements and by coefficients preserves that span.
Consequently, in the fixed basis of \(A\), membership in
this extended ideal means membership of every
\(B[G]\)-coordinate in the old ordinary ideal.

The two containments in (BCN.37a) give
\(\mathfrak m_B^N A\subset\mathfrak m_A^N\) and
\(\mathfrak m_A^{cN}\subset\mathfrak m_B^N A\).
Thus every extended old closed relation belongs to each
new defining ideal at depth \(N\).  Conversely, a new
closed relation belongs at depth \(cN\), with the same
\(H\), to the extended old defining ideal at depth \(N\).
Its coordinates therefore belong to every old defining
ideal, and hence to \(I_B\).  There are only finitely many
coordinates of the fixed element, so they assemble into
\(AI_B\).  This proves (BCN.38).  The subgroup \(H\) is
unchanged throughout both depth comparisons.
Our actual \(A_B\) satisfies this hypothesis with \(c=1\)
by BCN3.  The general statement also applies to a finite
free local splitting algebra with merely cofinal maximal
ideal topologies.
\hfill\(\square\)
% END BCN17

% BEGIN BCN18
\paragraph{BCN18. The exact closed ideal after adjoining \(s\).}\label{claim:BCN18}
Put \(T_0=A[[s]]=\mathscr T(B)\).
Extend the determinant and \(Z\) by the constant map.
The closed ideal with respect to
\(\mathfrak m_{T_0}=(\mathfrak m_A,s)\) is
\[
 I_{T_0}=T_0\,I_A=T_0\,I_B
 \quad\hbox{inside }T_0[G].                                    \tag{BCN.39}
\]
This is an ordinary generated ideal, not an assertion that
every infinite-support series belongs to the group algebra.

\emph{Proof.}
Let \(z\in I_{T_0}\) and fix its finite group support \(S\).
It has the expansion
\[
 z=\sum_{r\ge0} z_rs^r,\qquad z_r\in A[S]\subset A[G].
                                                               \tag{BCN.40}
\]
Fix \(r,N\) and an original open normal subgroup \(H\).
Use the defining membership for \(z\) at the single depth
\(N+r\) and this same \(H\).  By BCN16 the algebraic
Cayley--Hamilton and \(Z\) part of the defining ideal is
generated by constants from its \(A[G]\) counterpart;
the coefficient of order \(r\) of any product involving
one such generator still lies in the corresponding
ordinary ideal of \(A[G]\).
The same is true of products involving \([h]-1\), for
\(h\in H\); the subgroup has not been changed.

The remaining generators have coefficients in
\((\mathfrak m_A,s)^{N+r}\).  This power is
\[
 \sum_{j=0}^{N+r}s^j\mathfrak m_A^{\,N+r-j}T_0.
\]
Only \(j\le r\) can contribute to the coefficient of
\(s^r\), and each contribution belongs to
\(\mathfrak m_A^{\,N+r-j}\subset\mathfrak m_A^N\).
Multiplying by arbitrary group-algebra elements does
not weaken that containment.  Extracting the coefficient
in the defining finite ideal expression therefore gives
\[
 z_r\in
 \left\langle\mathrm{CH}(\Delta_A),Z_A,\mathfrak m_A^N[1],
                  [h]-1\ (h\in H)\right\rangle .
\]
Since \(N,H\) were arbitrary, \(z_r\in I_A\).

Now use the fixed support \(S\): the submodule
\(I_A\cap A[S]\) of the finite free \(A\)-module \(A[S]\)
is finitely generated because \(A\) is Noetherian.
Choose generators \(w_1,\ldots,w_b\) in that intersection.
For each \(r\), choose \(a_{rj}\in A\) with
\(z_r=\sum_j a_{rj}w_j\).  Equation (BCN.40) then becomes
\[
 z=\sum_{j=1}^b\left(\sum_{r\ge0}a_{rj}s^r\right)w_j.
                                                               \tag{BCN.41}
\]
This is a finite sum in \(T_0[G]\), so \(z\in T_0I_A\).
Each \(w_j\) lies in \(AI_B\) by BCN17, proving the
forward containment in (BCN.39).

Conversely, for \(w\in I_A\), apply its defining
approximation at depth \(N\) and subgroup \(H\), then
include its coefficients as constant series.  Since
\(\mathfrak m_A^N\subset\mathfrak m_{T_0}^N\), its image
belongs to every defining ideal for \(I_{T_0}\).
As \(I_{T_0}\) is a two-sided ideal it contains \(T_0I_A\).
This proves equality and completes all simultaneous
adic and group-quotient bookkeeping.
\hfill\(\square\)
% END BCN18

% BEGIN BCN19
\paragraph{BCN19. The literal augmentation ideal of the new pair.}\label{claim:BCN19}
For the pair \(u'\) in BCN7, let
\(I=\mathrm{augmentationIdealKernel}(u)\), interpreted
as the underlying set of its closed two-sided ideal.
Then
\[
 \left\langle\mathrm{augmentationIdealKernel}(u')\right\rangle
 =
 \left\langle\iota_B(I)\right\rangle
 \quad\hbox{in }\mathscr T(B)[G].                               \tag{BCN.42}
\]
In fact the left underlying closed ideal itself is the
right generated ideal.

\emph{Proof.}
The original \(\mathrm{pairGens}\) consists of the
ordered ordinary elements \(oEl(v,\psi,\gamma)\), their
polarized \(\star\)-images, the elements \([g]-1\) for
\(g\in r.I_p(w)\), and their polarized images.
BCN7 gives equality of all \(I_p(w)\)'s, and the character
values are coefficient images.  The definition of
\(\star\) uses the same cyclotomic coefficient map
\(\mathbb Z_\ell\to O\to O'\); it therefore commutes with
\(\iota_B\).  Thus the new generating set is precisely
the coefficient image of the old generating set, with
all group indices and orderings retained.
Apply BCN18 to this set and \(\Delta=u.D\).
The source's augmentation ideal is already the set of
all elements of this closed ideal, so taking its ordinary
two-sided span makes no further change.

This proof is specific to the actual unramified extension
constructed in BCN2 and the literal \(\mathrm{Til}\) ring.
It does not assert the source's more general
\(\mathrm{closure\_bc}\) for arbitrary local homomorphisms.
It requires no unreconstructed \(\mathrm{finE}\),
\(\mathrm{pair\_mf}\), or finite-dimensional
Cayley--Hamilton quotient premise.
\hfill\(\square\)
% END BCN19

% BEGIN BCN20
\paragraph{BCN20. Semisimplicity descends under a field extension
without a finite-dimensional hypothesis.}\label{claim:BCN20}
Let \(E\) be any algebra over a field \(\kappa\), and let
\(K/\kappa\) be any field extension.  If \(K\otimes_\kappa E\)
is semisimple as a ring, then so is \(E\).

\emph{Proof.}
Use the source's actual convention: a ring is semisimple
when every submodule of its left regular module has a
complement.  Equivalently, every left ideal has a
left-module retraction.  This is the definition in the
mathlib commit pinned by the source, not a convention
that builds finite dimension into the term.

Choose a \(\kappa\)-linear functional \(b:K\to\kappa\)
with \(b(1)=1\), by extending \(1\) to a vector-space basis.
For a left ideal \(L\subset E\), the tensor
\(K\otimes_\kappa L\) embeds in \(K\otimes_\kappa E\),
since vector-space extension is exact.
It is a left ideal, and hence has a
\(K\otimes_\kappa E\)-linear retraction \(p\).
The composite
\[
 E\xrightarrow{\,e\mapsto1\otimes e\,}
 K\otimes E\xrightarrow{p}K\otimes L
 \xrightarrow{\,b\otimes1\,}L
                                                               \tag{BCN.43}
\]
is \(E\)-linear: coefficients from \(K\) commute with
the \(E\)-action.  On \(L\) it is the identity because
\(p(1\otimes l)=1\otimes l\) and \(b(1)=1\).
Thus \(E=L\oplus\ker((\mathrm{BCN.43}))\) as left modules.
Every left ideal splits, proving the exact semisimple-ring
predicate.  No finite-dimensionality, separability, or
characteristic restriction was used.
\hfill\(\square\)
% END BCN20

% BEGIN BCN21
\paragraph{BCN21. Generic same-prime transport of nonsemisimplicity.}\label{claim:BCN21}
Let \(f:B\to S\) be any homomorphism of commutative rings,
let \(\mathfrak q'\subset S\) be prime and
\(\mathfrak q=f^{-1}\mathfrak q'\).
For sets \(I\subset B[G]\) and \(I'\subset S[G]\), suppose
\[
 I'\subset\langle f(I)\rangle .                               \tag{BCN.44a}
\]
Then \(\mathrm{NonSS}(I,\mathrm{id}_B,\mathfrak q)\)
implies \(\mathrm{NonSS}(I',\mathrm{id}_S,\mathfrak q')\).
In particular, if \(\mathfrak q'\subset\mathscr T(B)\)
is prime, \(\mathfrak q=\iota_B^{-1}\mathfrak q'\), and
\[
 \mathrm{NonSS}
   (\mathrm{augmentationIdealKernel}(u),\mathrm{id}_B,
    \mathfrak q)
\]
holds, then the identical predicate for \(u'\) holds
at \(\mathfrak q'\).

\emph{Proof.}
Take a witnessing field \(\kappa\) and map \(b:B\to\kappa\)
with kernel \(\mathfrak q\), such that
\[
 E=\kappa[G]/\langle b(I)\rangle
 \quad\hbox{is not semisimple}.                                \tag{BCN.44}
\]
Let \(\kappa'=\operatorname{Frac}(S/\mathfrak q')\).
Both fields contain
\(F=\operatorname{Frac}(B/\mathfrak q)\), with the same
map from \(B\).  The tensor product
\(\kappa\otimes_F\kappa'\) is nonzero: a basis of one
field over \(F\) stays a basis after scalar extension.
Choose a prime of this commutative ring and pass to the
fraction field \(K\) of its quotient.  The induced maps
from \(\kappa,\kappa'\) into \(K\) are injective, because
they are unital maps of fields.  They agree on \(F\).

Use the map \(S\to\kappa'\to K\) as the new field witness;
its kernel is exactly \(\mathfrak q'\).
Group algebras commute with field extension by their
basis \([g]\), and quotienting by a generated two-sided
ideal commutes with tensoring.  Hypothesis (BCN.44a)
therefore gives a surjection from the new quotient onto
\(K\otimes_\kappa E\).  In our actual application BCN19
gives equality of the two ideals, and hence the stronger
identification
\[
 K[G]/\langle I'\rangle
 \simeq K\otimes_\kappa E .                                    \tag{BCN.45}
\]
If the new quotient were semisimple, so would be its
quotient \(K\otimes_\kappa E\): left ideals in a quotient
split by taking the image of a complement of their inverse
image.  BCN20 would then make \(E\) semisimple,
contradicting (BCN.44).
This is a field witness for the source's original
existential \(\mathrm{NonSS}\) definition, with no
finite-spanning auxiliary hypothesis.
\hfill\(\square\)
% END BCN21

% BEGIN BCN22
\paragraph{BCN22. Discriminant stability uses the same two tuples.}\label{claim:BCN22}
For every ideal \(\mathfrak q'\subset\mathscr T(B)\),
\[
 \mathrm{DiscSt}(\Delta,\iota_B^{-1}\mathfrak q')
 \quad\Longleftrightarrow\quad
 \mathrm{DiscSt}(u'.D,\mathfrak q').
                                                               \tag{BCN.46}
\]
Together with BCN21 this transports the original
\(\mathrm{StOrTh}\) predicate forward to every lifted prime.

\emph{Proof.}
The original \(\mathrm{DiscSt}\) asks for two
\(n^2\)-tuples \(g,h\) of group elements whose matrix
\[
 \big(\Delta.\operatorname{tr}([g_i h_j])\big)_{i,j}
\]
has determinant outside the prime.
For every such pair of tuples, (BCN.18) identifies the new matrix
with the coefficient image of the old one.
The determinant is a polynomial with integer coefficients,
so its new value is the image of its old value.
Membership in \(\mathfrak q'\) is exactly membership
of the old determinant in its contraction.
Both tuples are unchanged in both directions, proving
(BCN.46).  The alternative \(\mathrm{NonSS}\) in
\(\mathrm{StOrTh}\) is carried forward by BCN21.
\hfill\(\square\)
% END BCN22

\subsection{The complete implication and the selected one-point version}

% BEGIN BCN23
\paragraph{BCN23. The full original implication
\(\mathrm{Core}\Rightarrow\mathrm{CoreBC}\).}\label{claim:BCN23}
Assume the original \(\mathrm{Core}\).
Then the original \(\mathrm{CoreBC}\) holds, with every
coefficient ring, residual datum, pair, ideal \(J\),
integer \(D\), supplied empty-parameter \(\mathrm{TWAut}\),
and every finite point chain quantified exactly as in
\texttt{PatchedSupportCompletions.lean}.

\emph{Proof.}
Fix arbitrary data of \(\mathrm{CoreBC}\):
\(\ell\) prime and \(\ell\ne2\); the original
\(L,G_0,T,n,O,r,B,u,J,D,\mathcal A\);
\(T.\mathrm{ProL}\); adaptedness; and
\(\mathrm{GenBy}_O(\mathrm{pairVals}(u))\).
Fix any \(h\ge0\) and points
\[
 x_i:\mathrm{Pt}(T,u,D+nT.q),\qquad 0\le i\le h,
 \qquad |x_i.x.K|=|x_0.x.K| .                                  \tag{BCN.47}
\]
Choose the single \(O'\) from BCN1--BCN2 using \(K_0=x_0.x.K\).
BCN3, BCN7, and BCN11 give the new coefficient and
automorphic data \(u',J',\mathcal A'\).
The group \(T\) and therefore its \(\mathrm{ProL}\) field
are unchanged.  Adaptedness and the required dense
generation including \(s\) are BCN7.
These choices are made for the whole point chain,
before any of its intermediate prime ideals are considered.

To prove \(\mathrm{ChainConcl}\), now fix an arbitrary
list \(e_j\) satisfying its prime, \(\mathrm{NonSS}\),
and adjacency conditions, a prime \(q_1\) with
\[
 J\subset q_1\subset p_0,\qquad
 \mathrm{DiscSt}(\Delta,q_1),
\]
and a final prime \(q\subset p_h\) with
\(\mathrm{StOrTh}(u,q)\), where \(p_i=\ker f_i\).
These are exactly the universally quantified tests of
\(\mathrm{ChainConcl}\).
BCN14 supplies \(P_i,E_j\) over \(p_i,e_j\).
Going down within \(P_0,P_h\) supplies primes
\[
 Q_1\subset P_0,\quad Q_1\cap B=q_1,\qquad
 Q\subset P_h,\quad Q\cap B=q .
\]
For each \(P_i\), BCN12--BCN13 lift the original point
to \(x_i'\), keeping its entire \(\mathrm{NP}\) and
\(\mathrm{Kap}\) data fixed and using the same field \(K_i\).
Its coefficient map \(O'\to K_i\) is onto and the
literal added parameter maps to \(t_i\).
The lifted point maps may depend on these lifted primes;
this is permitted because \(\mathrm{Core}\) quantifies
over every such point chain.

Apply \(\mathrm{Core}\) to the new data and the full
chain \(x_i'\).
Use the intermediate primes
\(\widetilde E_j=E_j\mathscr T(B)\), the initial
prime \(\widetilde Q_1\), and the final prime
\(\widetilde Q\).
They are prime and have their required containments
in the kernels of \(f_i'\), by BCN15.
BCN21 transports every intermediate \(\mathrm{NonSS}\).
BCN22 transports the initial discriminant condition
and the final \(\mathrm{StOrTh}\).
The original containment \(J\subset q_1\), together
with contraction in (BCN.34), is exactly
\(J'\subset\widetilde Q_1\).
Thus the \(\mathrm{ChainConcl}\) supplied by
\(\mathrm{Core}\) gives
\[
 J'\subset\widetilde Q .
\]
For every \(b\in J\), its image \(\iota_B(b)\) belongs
to \(J'\), hence to \(\widetilde Q\); contraction in
(BCN.34) gives \(b\in q\).  Therefore \(J\subset q\),
which is the original desired conclusion.
All \(h\), all equal-cardinality chains, all intermediate
primes and both endpoint primes were arbitrary.
This proves the complete original \(\mathrm{CoreBC}\).
\hfill\(\square\)
% END BCN23

% BEGIN BCN24
\paragraph{BCN24. Only the selected normalized Core endpoint is needed
for the selected one-point route.}\label{claim:BCN24}
Restrict both source contracts to \(h=0\), to the original
full block \(x.x.a=n\), and to \(x.k.g=0\).
A supplied normalized \(\mathrm{Core}\) conclusion for
these data alone implies the equally restricted
\(\mathrm{CoreBC}\) conclusion.
It does not require the arbitrary-\(h\), arbitrary-\(g\)
antecedent of BCN23.

\emph{Proof.}
Run the construction of BCN23 for its single original
point \(x_0\).  There are no \(e_j\) and no intermediate
\(\mathrm{NonSS}\) tests.  There are still the initial
stable prime \(q_1\), the final \(\mathrm{StOrTh}\) prime
\(q\), and their lifts under one \(P_0\); BCN14--BCN15
supply them exactly as before.
By BCN13 the normalized point has
\[
 n'=n,\qquad x_0'.x.a=x_0.x.a=n,\qquad
 x_0'.k.g=x_0.k.g=0 .                                         \tag{BCN.48}
\]
Its \(\mathrm{B1}\), \(\mathrm{B2}\), \(D+nT.q\), group
packet, and same point field are unchanged.
Invoke only the supplied normalized one-point conclusion
at these data.  Its containment \(J'\subset\widetilde Q\)
contracts to \(J\subset q\), as before.
For the selected rank-four route, this gives \(n=a=4\)
and \(g=0\) throughout.  The coefficient normalization
therefore does not impose a new requirement to prove
a full arbitrary-chain Core theorem first.
\hfill\(\square\)
% END BCN24

% BEGIN BCN25
\paragraph{BCN25. Composition with the already reconstructed pro-\(\ell\)
descent.}\label{claim:BCN25}
The complete original \(\mathrm{Core}\) implies the complete
original \(\mathrm{Raw}\).
Likewise the selected normalized one-point full-block
Core endpoint suffices, pointwise, for its selected Raw
endpoint.  The original \(B,J\), final prime, point fields,
block sizes, \(\mathrm{Kap}\), and all Selmer-bound
quantifiers are preserved at the final conclusion.

\emph{Proof.}
BCN23 gives \(\mathrm{CoreBC}\).
Inherited PQL20 proves the original
\(\mathrm{CoreBC}\Rightarrow\mathrm{Raw}\) implication,
with the full quotient group, pair, point, module and
ideal transports already reconstructed there.
This proves the first assertion.

For the second assertion, use BCN24 in that same
argument.  The PQL13--PQL17 point construction keeps
the original field, representation block dimensions,
\(\mathrm{Kap}\) size, eigenvalues and the coupled
\(\mathrm{B1}/\mathrm{B2}\) witnesses, while replacing
only the necessary group presentation.
PQL20 invokes its \(\mathrm{CoreBC}\) consumer on just
this descended chain; for a single full-block point it
therefore needs only the restricted consumer provided
by BCN24.  Its final ideal contraction is the same
one-way \(\mathrm{NonSS}\) transport and literal
\(J\)-containment already proved in PQL18--PQL20.
Thus there is no hidden use of a general arbitrary-\(g,h\)
Core theorem in the selected composition.

Neither assertion constructs the antecedent normalized
patched support, any initial \(\mathrm{TWAut}\), the four
Hida endpoints on the same \(J\), or the arithmetic
ordinary/control/type data.  Those remain the original
upstream tasks.  The coefficient-normalization and
added-Unit-parameter implication itself has no additional
unproved supplier.
\hfill\(\square\)
% END BCN25

\subsection{Scope, source correspondence, and next use}

BCN1--BCN2 supply every clause of the actual
\texttt{PatchingBaseChange/Oprime.lean:oprime}.
BCN3--BCN11 supply the needed original
\texttt{Til}, residual, pair, base, ring, map, and
\texttt{aut\_bc} construction for that particular
finite free unramified extension.
BCN12--BCN15 give the exact point/kernel and prime-chain
interfaces of \texttt{pt\_lift}, \texttt{chain\_tens}
and \texttt{til\_prime}.
BCN16--BCN22 reconstruct the coefficient-extension and
stability steps in the source's
\texttt{AutomorphicDetection/\allowbreak{}CoreBC.lean} by a stronger
specialized route: the finite-support argument proves
the actual closed ideal identity, and the scalar-functional
argument proves the exact semisimplicity descent without
a finite-dimensional premise.
The source's general \texttt{finE} and general
\texttt{closure\_bc} theorems have not been reclassified
as independently reconstructed by this specialization.

The next same-object join is the original normalized
Core producer: its patched rings, maps and modules,
the exact support ideal and the normalized point's
local dimension/chain argument.  For the selected route,
BCN24 is the consumer to use with a supplied one-point
full-block producer.  It preserves \(D+nT.q\),
\(n=a=4\), \(g=0\), the one common coefficient ring for
the entire finite chain, the same \(J\) after contraction,
and every original quantifier in \(\mathrm{B1}\) and
\(\mathrm{B2}\).  No Lean replay or new kernel audit
is claimed.


